\begin{abstract}
Reacting to task-level failures in long-horizon manipulation is commonly formalized with a hybrid automaton whose discrete modes select the active skill; reactivity then hinges on a \emph{localization function} that grounds the continuous state to the correct mode; equivalently, on knowing where the \emph{guards} between modes lie. Prior work either hand-engineers this grounding, inheriting every error of the designer's analytic model, or learns it from perturbation-and-reset rollouts that must induce the very failures a deployed system exists to avoid. We propose SEAL (Safe, Efficient Active Learning), which learns this grounding of the task automaton from a handful of safe, human-segmented demonstrations: an LLM proposes a minimal, spatially invariant feature space, verified by an information-theoretic observability check; a Gaussian Process Classifier localizes the mode, its epistemic uncertainty directing additional safe demonstration requests to the states where the guard estimate is weakest; and its predictive confidence decouples switching latency from false-positive rate at runtime. No failure is ever induced during training. Across five simulated and real-world contact-rich tasks and a 10-participant user study, SEAL matches the accuracy of perturbation-based grounding and expert-tuned automata while incurring fewer false-positive transitions at equal latency, and it does so without resets, autonomous rollouts, or failure data.
\end{abstract}

\begin{IEEEkeywords}
Learning from Demonstration, Failure Detection and Recovery, Formal Methods in Robotics and Automation, Hybrid Logical/Dynamical Planning and Verification, Human-in-the-Loop.
\end{IEEEkeywords}
